%\documentclass[12pt,letterpaper]{article}



%%%
%% Required packages:
%%
\usepackage{etex}
\usepackage{amsmath}
\usepackage{amssymb}
\usepackage{amstext}
\usepackage{amsthm}
\usepackage{fancyhdr,fancybox}
\usepackage{mathrsfs} % need for \mathscr command
\usepackage{multicol}
\usepackage[normalem]{ulem}
%\usepackage{pictex,pictexplus}
% \usepackage{pictexwd}
\usepackage{rotating}
\usepackage{url}
\usepackage{xfrac}
\usepackage{booktabs}
\usepackage{color,soul}
\usepackage{comment}

\usepackage{hyperref}
\hypersetup{
    colorlinks=true,     % false: boxed links; true: colored links
    linkcolor=blue,      % color of internal links
    citecolor=blue,      % color of links to bibliography
    filecolor=blue,      % color of file links
    urlcolor=blue        % color of external links
}


\usepackage{tikz}
\usetikzlibrary{backgrounds,calc}

%%
%% Common formatting commands
%%
\setlength{\parindent}{0in}
\setlength{\textwidth}{7in}
\setlength{\evensidemargin}{-0.25in}
\setlength{\oddsidemargin}{-0.25in}
\setlength{\parskip}{.5\baselineskip}
\setlength{\topmargin}{-0.5in}
\setlength{\textheight}{9in}

%               Problem and Part
\newcounter{problemnumber}
\newcounter{partnumber}
\newcounter{subpartnumber}
\newcommand{\Problem}{%
  \refstepcounter{problemnumber}%
  \setcounter{partnumber}{0}%
  \setcounter{subpartnumber}{0}%
  \item[\makebox{\hfil\textbf{\theproblemnumber.}\hfil}]%
  }
\newcommand{\InProblem}[2][1.6]{%
  \refstepcounter{problemnumber}%
  \setcounter{partnumber}{0}%
  \setcounter{subpartnumber}{0}%
  \textbf{\theproblemnumber.}\ \ \parbox[t]{#1in}{#2}%
}
\newcommand{\InSmallProblem}[1]{\InProblem[1.05]{#1}}
\newcommand{\NoProblem}[1]{%
  \mbox{\hfil\textbf{#1}\hfil}%
  }
\newcommand{\UnProblem}{%
  \refstepcounter{problemnumber}%
  \setcounter{partnumber}{0}%
  \setcounter{subpartnumber}{0}%
  \mbox{\hfil\textbf{\theproblemnumber.}\hfil}%
  }
\newcommand{\InFlexProblem}[1]{%
  \refstepcounter{problemnumber}%
  \setcounter{partnumber}{0}%
  \setcounter{subpartnumber}{0}%
  \textbf{\theproblemnumber.}\ \ #1%
}
%%  Part:
\newcommand{\Part}{%
  \renewcommand{\thepartnumber}{(\alph{partnumber})}%
  \refstepcounter{partnumber}
  \setcounter{subpartnumber}{0}%
  \item[(\alph{partnumber})]%
}
\newcommand{\InPart}[2][1.75]{%
  \renewcommand{\thepartnumber}{(\alph{partnumber})}%
  \refstepcounter{partnumber}%
  \setcounter{subpartnumber}{0}%
  (\alph{partnumber})\ \ \parbox[t]{#1in}{#2}%
}
\newcommand{\UnPart}{%
  \refstepcounter{partnumber}%
  \renewcommand{\thepartnumber}{(\alph{partnumber})}%
  \setcounter{subpartnumber}{0}%
  (\alph{partnumber})%
}
\newcommand{\InLargePart}[1]{\InPart[3.05]{#1}}
\newcommand{\InFullPart}[1]{\InPart[6.2]{#1}}
\newcommand{\InSmallPart}[1]{\InPart[1.05]{#1}}
%% SubPart:
\newcommand{\SubPart}{%
  \refstepcounter{subpartnumber}%
  \item[(\roman{subpartnumber})]%
  }
\newcommand{\InSubPart}[2][2.25]{%
  \stepcounter{subpartnumber}%
  (\roman{subpartnumber})\ \ \parbox[t]{#1in}{#2}%
}
\newcommand{\InSmallSubPart}[1]{\InSubPart[1.5]{#1}}
\newcommand{\InTinySubPart}[1]{%
  \stepcounter{subpartnumber}%
  (\roman{subpartnumber})\ \ \makebox{#1}%
}
\newcommand{\InMySubPart}[2][2.25]{%
  \stepcounter{subpartnumber}%
  \parbox[t]{#1in}{\makebox[0.3in][l]{(\roman{subpartnumber})}\ \ {#2}}%
}
\newcommand{\TrueFalse}{\textbf{T}\ \ \textbf{F}\ \ }
\newcommand{\TFPart}[1]{% 
  \stepcounter{partnumber}%
  \setcounter{subpartnumber}{0}%
  \item[(\alph{partnumber})] \TrueFalse\parbox[t]{5.5in}{#1}}

\newcommand{\Points}[1]{(#1 points) \ }
\newcommand{\Point}[1]{(#1 point) \ }


%% For Answers:
\newcounter{answernumber}
\newcommand{\Answer}{%
  \refstepcounter{answernumber}%
  \setcounter{partnumber}{0}%
  \setcounter{subpartnumber}{0}%
  \item[\makebox{\hfil\textbf{\theanswernumber.}\hfil}]%
  }


% Setting up the headers for the first page
%\chead{} % will default to what is typed in the file.
\fancyhead[L]{\textsc{Math 22a}}
\fancyhead[R]{\textsc{Fall 2024}}
\fancyfoot[R]{
	\vspace*{\fill}\footnotesize\textcopyright{} \emph{Harvard Math 22a}	
}
\renewcommand{\headrulewidth}{0pt}

%\fancyhead[C]{}
%	\chead{}	
%	\lhead{}
%	\rhead{}
%	\cfoot{}


% Putting copyright on the footer of each page
%\fancypagestyle{secondpage}{
%	\fancyhead[C]{TESTing again} % change this to be blank
%		%\chead{}	
%	\fancyhead[L]{bears} % change this to be blank
%	\fancyhead[R]{dogs} % change this to be blank
%	\fancyfoot[R]{ %keeps the footer the same
%		\vspace*{\fill}\footnotesize\textcopyright{} \emph{Harvard Math 22a}	
%	}
%}
%\renewcommand{\headrulewidth}{0pt}
%\pagestyle{fancy}
\pagestyle{empty}


%\lhead{}
%\rhead{}
%\rfoot{\vspace*{\fill}\footnotesize\textcopyright{} \emph{Harvard Math 22a}}
%\renewcommand{\headrulewidth}{0pt}
%\pagestyle{fancy}




%%
%% Common shortcut commands:
%%
\newcommand{\ds}{\displaystyle}
\newcommand{\sgn}{\operatorname{sgn}}
\renewcommand{\Pr}{\operatorname{P}}
\newcommand{\CPr}[3][{}]{\Pr_{\text{#1}}\left(#2\,|\,#3\right)}

\newcommand{\C}{\mathbb{C}}
\newcommand{\R}{\mathbb{R}}
\newcommand{\Z}{\mathbb{Z}}
\newcommand{\N}{\mathbb{N}}
\newcommand{\Q}{\mathbb{Q}}
\newcommand{\D}{\mathbb{D}}

\newcommand{\rref}{\operatorname{rref}}
\newcommand{\im}{\operatorname{im}}
\newcommand{\rank}{\operatorname{rank}}
\newcommand{\diag}{\operatorname{diag}}
\newcommand{\Span}{\operatorname{span}}
%\renewcommand{\vec}[1]{\mathbf{#1}}
\newcommand{\charpoly}{\mathscr{P}}
\newcommand{\nicefrac}[2]{\sfrac{#1}{#2}} % using xfrac package
\newcommand{\topology}{\mathcal{T}}
\newcommand{\Inn}{\operatorname{Inn}}

\newcommand{\blankbox}[2]{\fbox{\rule{#1}{0in}\rule{0in}{#2}}}

\newenvironment{myindentpar}[1]%
 {\begin{list}{}%
         {\setlength{\leftmargin}{#1}
          \setlength{\rightmargin}{\leftmargin}}%
         \item[]%
 }
 {\end{list}}
\newtheorem{theorem}{Theorem}
\newtheorem*{NStheorem}{Theorem}
\newtheorem{cor}[theorem]{Corollary}
\newtheorem*{claim}{Claim}
\newtheorem{defn}{Definition}

\newenvironment{amatrix}[1]{%
  \left(\begin{array}{@{}*{#1}{c}|c@{}}
}{%
  \end{array}\right)
}

%--------grstep
% For denoting a Gauss' reduction step.
% Use as: \grstep{\rho_1+\rho_3} or \grstep[2\rho_5 \\ 3\rho_6]{\rho_1+\rho_3}
\newcommand{\grstep}[2][\relax]{%
   \ensuremath{\mathrel{
       {\mathop{\longrightarrow}\limits^{#2\mathstrut}_{
                                     \begin{subarray}{l} #1 \end{subarray}}}}}}
\newcommand{\swap}{\leftrightarrow}




\section{{Matrix Equations and Solutions Sets,  $\vec\S$1.4}}% 
\chead{\Large\textbf{Matrix Equations and Solutions Sets,  $\vec\S$1.4}}% 

%\begin{document}
\thispagestyle{fancy}

Recall from last class:

\begin{itemize}

\item A matrix is in {\bf echelon form} if the following 3 properties hold:
\begin{enumerate}
\item All non-zero rows are above zero rows.
\item Each leading entry of a row is in a column to the right of the leading entry of the row above it.
%\item All entries in a column below a leading entry are 0.
\end{enumerate}
\item A matrix in echelon form is in {\bf reduced echelon form} if 
\begin{enumerate}
\item The leading entries are all 1.
\item Each leading 1 is the only non-zero entry in its column.
\end{enumerate}
\item Each matrix is row equivalent to one and only one reduced echelon matrix.
\item A {\bf pivot position} in a matrix $A$ is a location in $A$ that corresponds to a leading 1 in the reduced row echelon form. A {\bf pivot column} is a column that has a pivot position. 
\item {\bf Row Reduction Algorithm (Gaussian Elimination).}
\begin{enumerate}
\item The leftmost non-zero column is a pivot column.
\item Select a non-zero entry in the pivot column as pivot. Interchange rows to move this entry to the top.
\item Use elementary row operations to create zeros below the pivot.
\item Apply 1-3 to the submatrix obtained by removing pivot row and column. Repeat until all non-zero rows are modified.
\item Beginning with rightmost pivot and working upwards and to the left, create 0's above each pivot. Scale pivot positions entries to be 1.
\end{enumerate}
\item The variables corresponding to the pivot columns in an augmented matrix are called {\bf basic} and the others are called {\bf free}.
\item A linear system is \textbf{consistent} if and only if the rightmost column in the augmented matrix is not a pivot column. If the system is consistent, then it has a unique solution when there are no free variables and infinite solutions when there is one or more free variables.
\item A matrix with 1 column is called a {\bf column vector}. A matrix with one row is called a {\bf row vector}.
\item We can add vectors of the same size by adding the corresponding components; namely,
if $\vec{u}=\begin{bmatrix} u_1\\ \vdots\\ u_n \end{bmatrix}$ and  $\vec{v}=\begin{bmatrix} v_1\\ \vdots\\ v_n \end{bmatrix}$, then $\vec{u}+\vec{v} =\begin{bmatrix} u_1+v_1\\ \vdots\\ u_n+v_n \end{bmatrix}$. Similarly, we can scale a vector by a real number $c$, by multiplying all the components by $c$; namely, $c\vec{u}=\begin{bmatrix} cu_1\\ \vdots\\ cu_n \end{bmatrix}$.
\item Given $\vec{v}_1,\dots, \vec{v}_p \in \mathbb{R}^n$ and scalars $c_1, \dots c_p$, the vector $$\vec{y} = c_1 \vec{v}_1 +\dots c_p \vec{v}_p$$ is called the {\bf linear combination} of $\vec{v}_1,\dots, \vec{v}_p$ with weights $c_1,\dots,c_p$.
\item A vector equation $$x_1 \vec{a}_1 + x_2 \vec{a}_2 +\dots + x_n \vec{a}_n =\vec{b}$$ has the same solutions as the augment matrix $$[\vec{a}_1 \; \vec{a}_2 \; \dots \; \vec{a}_n \Big| \vec{b}].$$ Hence $\vec{b}$ is a linear combination of $\vec{a}_1, \dots, \vec{a}_n$ if and only if there is a solution to this matrix.
\item If $\vec{v}_1, \dots, \vec{v}_p \in \mathbb{R}^n$, then the set of all linear combinations of $\vec{v}_1, \dots, \vec{v}_p$ is called the {\bf span} of $\vec{v}_1, \dots, \vec{v}_p$. We write this set as $\text{Span} \{\vec{v}_1, \dots, \vec{v}_p\}$.
\end{itemize}

\newpage

Let $A$ be an $m \times n$ matrix, then $$A = [ \vec{a}_1\; \dots \; \vec{a}_n], \; \; \vec{a}_j \in \mathbb{R}^m,\; 1\leq j \leq n.$$  
If $x \in \mathbb{R}^n$, then we define $$A \vec{x} = \vec{a}_1 x_1 + \vec{a}_2 x_2 +\dots + \vec{a}_n x_n.$$ Thus the matrix equation $A\vec{x} =\vec{b}$ is the same as the vector equation $$\vec{a}_1 x_1 + \vec{a}_2 x_2 +\dots + \vec{a}_n x_n=\vec{b}.$$


\begin{enumerate}

\Problem Given $$A=\begin{bmatrix}
1 & 3 & -4\\
1 & 5 & 2\\
-3 & -7 & 6\\
\end{bmatrix}\text{  and }\vec{b}=\begin{bmatrix} -2\\ 4 \\ 12\\ \end{bmatrix}.$$ Is $A\vec{x} =\vec{b}$ consistent?

\newpage

{\bf Fact:} $A\vec{x} =\vec{b}$ has a solution if and only if $\vec{b} \in \text{Span}\{\vec{a}_1,\dots, \vec{a}_n\}$ where $\vec{a}_j$ is the $j$-th column of $A$.

\Problem Let $A=\begin{bmatrix} 1 & 2 & 3 \\ 4 & 5 & 6\\ 7 & 8 & 9\\ \end{bmatrix}$. Is $A\vec{x} =\vec{b}$ consistent for all $\vec{b} \in \mathbb{R}^3$.

\vfill

\begin{defn}
A set of vectors $\vec{v}_1, \dots, \vec{v}_p$ in $\mathbb{R}^m$ {\bf spans} $\mathbb{R}^m$ if every vector in $\mathbb{R}^m$ is a linear combination of $\vec{v}_1, \dots, \vec{v}_p$.
\end{defn}

\newpage

{\bf Theorem} Let $A$ be an $m \times n$ matrix, then the following are equivalent.
\begin{enumerate}
\item For each $\vec{b} \in \mathbb{R}^m$, the matrix equation $A \vec{x} =\vec{b}$ is consistent.
\item Each $\vec{b} \in \mathbb{R}^m$ is a linear combination of the columns of $A$.
\item The columns of $A$ span $\mathbb{R}^m$.
\item $A$ has a pivot position in every row.
\end{enumerate}

\Problem Prove the theorem.

\newpage

{\bf Fact (or Theorem 5 from the book):} If $A$ is a $m \times n$ matrix and $\vec{u}, \vec{v} \in \mathbb{R}^n$, $c\in \mathbb{R}$, then
\begin{enumerate}
\item $A(\vec{u} +\vec{v})=A\vec{u} +A\vec{v}$
\item $A(c \vec{u}) =c (A\vec{u}).$
\end{enumerate}

\Problem Prove the result.

\vfill

\Problem {\bf True or False}. If the augmented matrix has a pivot position in each row, then it is consistent.

\vfill

\Problem Let $$\vec{u} =\begin{bmatrix} 0 \\ 4 \\ 4\\ \end{bmatrix} \text{ and } A=\begin{bmatrix} 2 & -5\\ 2 & 6\\ 1 & 1\\ \end{bmatrix}.$$ Is $\vec{u}$ in the span of the columns of $A$?

\vfill


\end{enumerate}

%%solutions start below
%\end{document}
\begin{comment}

\newpage
\chead{\Large\textbf{Matrix Equations and Solutions -- Answers / Solutions}}%
\lhead{}
\rhead{}
\thispagestyle{fancy}

\begin{enumerate}
  \Answer
  An echelon form for $[A|b]$ is:
  \[
    \left[
      \begin{array}{ccc|c}
        1&3&-4&-2\\
        0&2&6&6\\
        0&0&-12&0
      \end{array}
    \right]
  \] and the last column is not a pivot, so the system is consistent.

  \Answer An echelon form for $[A|b]$ is:
  \[
    \left[
      \begin{array}{ccc|c}
        1&2&3&b_1\\
        0&-3&-6&b_2-4b_1\\
        0&0&0&b_3-2b_2+b_1
      \end{array}
    \right]
  \] where $b=\left[\begin{array}{c} b_1\\b_2\\b_3\end{array}\right]\in \mathbb{R}^3$ is an arbitrary vector. The last column is a pivot if $b_3-2b_2+b_1\neq 0$, so the system is consistent if $b_3-2b_2+b_1=0$. 

  \Answer
  (a)$\implies$(b): Suppose that (a) is true, and let $\vec{b}\in \mathbb{R}^m$. Then since $A\vec{x}=\vec{b}$ is consistent, there is at least one solution $x=\left[\begin{array}{c}x_1\\x_2\\ \vdots\\ x_n\end{array}\right]$. Let $\vec{a_i}$ be the $i^{th}$ column of $A$. Then by definition of matrix multiplication, $\vec{b}=\sum_{i=1}^n x_i\vec{a_i}$. Since $\vec{b}$ was arbitrary, any vector in $\mathbb{R}^m$ is a linear combination of the columns of $A$. 
  
  (b)$\implies$(c): Suppose that (b) is true. In order for the columns of $A$ to span $\mathbb{R}^m$, then any $\vec{b}\in \mathbb{R}^m$ must be a linear combination of the columns of $A$. This is exactly what (b) says.

  (c)$\implies$(d): Suppose that (d) is false, i.e. that $A$ does not have a pivot position in every row, and let $\vec{b}\in \mathbb{R}^m$ be an arbitrary vector. Then by the row reduction algorithm, $[A|\vec{b}]$ is row equivalent to an augmented matrix $[B|\vec{c}]$ where $B$ has at least one zero row at the bottom.

  If the last entry of $\vec{c}$ is not zero, then the last column of $[B|\vec{c}]$ will be a pivot, so will not be consistent. Since row equivalent matrices have the same solution set, then $[A|\vec{b}]$ is also inconsistent, so there is no solution to $A\vec{x}=\vec{b}$, so $\vec{b}$ is not in the span of the columns of $A$.

  Since row operations are reversible, just pick a $\vec{c}$ with the last entry nonzero, and then undo the row operations to get $\vec{b}$. 

  (d)$\implies$(a): Suppose (d) is true, and let $\vec{b}\in \mathbb{R}^m$ be arbitrary. Form the augmented system $[A|\vec{b}]$ and perform row reduction. By the assumption that $A$ has a pivot position in every row, then we will never encounter the situation where the last column is a pivot column. This is because to get the last column as a pivot column, we need to be able to perform row operations to get a zero row in $A$, which cannot happen by our assumptions. Thus the system $A\vec{x}=\vec{b}$ will always be consistent.

  \Answer
  \begin{enumerate}
  \item Let $A,\vec{u},\vec{v}$ be as follows:
    \[
      A = \left[
        \vec{a}_1 \; \cdots \; \vec{a}_n
      \right]\quad \vec{u}=
      \left[
        \begin{array}{c}
          u_1\\
          u_2\\
          \vdots\\
          u_n
        \end{array}
      \right]\quad \vec{v}=
      \left[
        \begin{array}{c}
          v_1\\
          v_2\\
          \vdots\\
          v_n
        \end{array}
      \right]
    \]

    Then we compute using the definitions:
\begin{align*}
A(\vec{u}+\vec{v})= A \begin{bmatrix}
u_1+v_1\\
u_2+v_2\\
\vdots\\
u_n+v_n\\
\end{bmatrix}&= (u_1+v_1)\vec{a}_1+\dots+(u_n+v_n)\vec{a}_n\\
&= u_1\vec{a}_1+v_1\vec{a}_1+\dots+u_n\vec{a}_n+v_n\vec{a}_n\\
&= (u_1\vec{a}_1+\dots+u_n\vec{a}_n)+(v_1\vec{a}_1+\dots+v_n\vec{a}_n)\\
&=A\vec{u}+A\vec{v}.
\end{align*}
    
  \item
   Exercise.
  \end{enumerate}

  \Answer This is {\bf false}. Consider the augmented matrix:
  \[
    \left[
      \begin{array}{cc|c}
        1&23&0\\
        0&0&1
      \end{array}
    \right]
  \]
  There is a pivot in each row, but the matrix is not consistent. A true statement here is that if augmented matrix does not have the last column as a pivot column, then it is consistent. 

  \Answer Remember that $\vec{u}$ being in the span of the columns of $A$ is equivalent to the augmented matrix $[A|\vec{u}]$ being consistent. An echelon form for $[A|\vec{u}]$ is:
  \[
    \left[
      \begin{array}{cc|c}
        2&-5&0\\
        0&11&4\\
        0&0&\frac{30}{11}
      \end{array}
    \right]
  \] so this system is not consistent, so $\vec{u}$ is not in the span of the columns of $A$. 

 \end{enumerate}
\end{comment}
%\end{document}


%%% Local ables: 
%%% mode: latex
%%% TeX-master: t
%%% End: 
